\section{Proposed CILP methodology}
\label{sec:methods}

\subsection{Problem formulation}
Given an attributed undirected graph $G=(V,E,X,Y)$ and edge-removal rate $r\in[0,1]$, return a sparsified edge set $E_r\subseteq E$ with
\[
|E_r|=(1-r)|E|
\]
(all nodes retained). Write $G_r=(V,E_r,X,Y)$ for the sparsified graph.
Downstream GCN Macro-F1 is the primary utility metric~\citep{kipf2017semi}.

\subsection{Advance relative to ILP-GCN}
Table~\ref{tab:advance} summarizes the methodological differences.
The social-network domain alone is not presented as novelty; the new element under study is composite counterfactual supervision with a matched task-only control under exact budgets.

\begin{table}[t]
\centering
\caption{\textbf{Methodological advance of CILP relative to ILP-GCN.}}
\label{tab:advance}
\scriptsize
\begin{tabular}{@{}p{2.2cm}p{5.3cm}p{5.8cm}@{}}
\toprule
Aspect & ILP-GCN~\citep{bangian2025ilp} & CILP \\
\midrule
Graph domain & Cheminformatics similarity graph & Attributed social graph \\
Edge target & Inverted link probability & Composite multi-criteria teacher \\
Criteria & Link score (and domain priors in the parent setting) & Task loss, connectivity, community-boundary proxy, degree-based spectral proxy, representation shift, group impact \\
Selection & Threshold / modularity-oriented sparsification in the parent paper & Exact-budget ranking by predicted importance $\mu_e$ \\
Uncertainty & Not used for ranking & Heteroscedastic $\sigma_e$ estimated but unused for ranking in reported runs \\
Evaluation & Domain-graph sparsification & Multi-seed social-network utility and structure under matched budgets \\
Main control & Random / classical baselines & Architecturally matched task-only teacher \\
\bottomrule
\end{tabular}
\end{table}

\subsection{Composite multi-criteria teacher}
CILP uses a composite multi-criteria teacher that combines deletion-induced effects with structural proxies (Fig.~\ref{fig:framework}).
Sample $n$ edges, delete each edge as required, and compute six normalized deletion-importance criteria (Supplementary Information~S5--S6):
\begin{itemize}[leftmargin=1.4em]
\item \textbf{Deletion-induced:} task loss, connectivity change, representation shift, group impact;
\item \textbf{Structural proxies:} community-boundary proxy, degree-based spectral proxy.
\end{itemize}
\[
y_e^{\mathrm{cf}}=\mathrm{minmax}\Big(\sum_k\beta_k\,\mathrm{minmax}(\Delta_k(e))\Big),
\]
with fixed weights $\beta=(1.0,0.5,1.0,0.3,0.5,0.5)$ for
(task, community-boundary proxy, connectivity, degree-based spectral proxy, representation, group),
identical across datasets and not development-tuned.
This is weighted-sum scalarization, not Pareto optimization.
The matched task-only teacher uses $y=\mathrm{minmax}(\Delta_{\mathrm{task}})$ with the same encoder and concatenation fusion.

\begin{figure}[t]
\centering
\resizebox{\textwidth}{!}{% Figure 1 — CILP framework
\begin{tikzpicture}[
  font=\sffamily\scriptsize,
  >=Latex,
  box/.style={draw=black!65, rounded corners=1.2pt, align=center, inner sep=2.6pt, fill=#1, minimum height=0.68cm},
  arr/.style={->, thick, black!60},
  dasharr/.style={->, thick, dashed, black!55},
  plab/.style={font=\sffamily\footnotesize\bfseries, anchor=west},
]
\definecolor{cA}{HTML}{EAF2F8}
\definecolor{cB}{HTML}{E8F6EF}
\definecolor{cC}{HTML}{FEF5E7}
\definecolor{cD}{HTML}{F5EEF8}
\definecolor{cE}{HTML}{F2F3F4}

\node[plab] at (-0.15,3.35) {(a) Training};
\node[box=cE, minimum width=1.45cm] (g) at (0.65,2.35) {Attributed\\social graph};
\node[box=cE, minimum width=1.55cm] (pp) at (2.5,2.35) {Preprocess\\+ splits};
\node[box=cA, minimum width=1.25cm] (enc) at (4.3,2.35) {Node\\encoder};
\node[box=cA, minimum width=1.45cm] (nv) at (6.05,2.85) {Node-centric\\edge view};
\node[box=cA, minimum width=1.45cm] (ev) at (6.05,1.85) {Edge-centric\\view};
\node[box=cD, minimum width=1.35cm] (fus) at (8.0,2.35) {Concat\\fusion};
\node[box=cB, minimum width=1.7cm] (dec) at (9.95,2.35) {Predicted\\importance $\mu_e$\\Est.\ uncertainty $\sigma_e$};
\node[box=cE, minimum width=1.5cm] (loss) at (11.9,2.35) {Training\\loss};

\draw[arr] (g)--(pp); \draw[arr] (pp)--(enc);
\draw[arr] (enc)--(nv); \draw[arr] (enc)--(ev);
\draw[arr] (nv)--(fus); \draw[arr] (ev)--(fus);
\draw[arr] (fus)--(dec); \draw[arr] (dec)--(loss);

\node[box=cC, minimum width=1.85cm] (samp) at (2.5,0.85) {Multi-criteria\\teacher};
\node[box=cC, minimum width=2.15cm] (del) at (4.85,0.85) {Six criteria:\\deletion effects\\+ structural proxies};
\node[box=cC, minimum width=1.45cm] (agg) at (7.15,0.85) {Normalize\\+ weight};
\node[box=cC, minimum width=1.25cm] (ycf) at (8.95,0.85) {Target\\$y^{\mathrm{cf}}_e$};
\draw[arr] (pp)--(samp); \draw[arr] (samp)--(del); \draw[arr] (del)--(agg); \draw[arr] (agg)--(ycf);
\draw[dasharr] (ycf.east) -- ++(0.35,0) |- (loss.south);
\node[font=\sffamily\tiny, text=black!60] at (10.55,1.35) {regression / ranking loss};

\node[plab] at (-0.15,-0.25) {(b) Exact-budget edge pruning};
\node[box=cB, minimum width=1.45cm] (sc) at (0.7,-1.15) {Trained\\scorer};
\node[box=cE, minimum width=1.45cm] (score) at (2.5,-1.15) {Score by\\$\mu_e$};
\node[box=cE, minimum width=1.35cm] (rank) at (4.3,-1.15) {Rank\\edges};
\node[box=cC, minimum width=1.85cm] (pr) at (6.4,-1.15) {Exact-budget\\edge pruning};
\node[box=cA, minimum width=1.45cm] (sp) at (8.4,-1.15) {Sparsified\\$G_r$};
\node[box=cD, minimum width=2.0cm] (evl) at (10.55,-1.15) {Downstream\\node classification};
\draw[arr] (sc)--(score); \draw[arr] (score)--(rank); \draw[arr] (rank)--(pr);
\draw[arr] (pr)--(sp); \draw[arr] (sp)--(evl);
\node[font=\sffamily\tiny, text=black!55, align=left] at (6.4,-1.9) {Deletion-induced: task, connectivity, representation, group.\\
Structural proxies: community, degree-based spectral.\\
Ranking uses $\mu_e$ only; $\sigma_e$ and constraints unused in reported runs.};
\end{tikzpicture}
}
\caption{\textbf{Schematic of Counterfactual Inverse Link Prediction (CILP).}
\textbf{a}~Forward scoring path with a parallel multi-criteria teacher providing targets $y^{\mathrm{cf}}_e$.
Four criteria are deletion-induced effects; community-boundary and degree-based spectral terms are structural proxies.
\textbf{b}~Exact-budget edge pruning by predicted importance $\mu_e$ and downstream node classification.
Reported runs rank by $\mu_e$ only; $\sigma_e$ and optional constraints are unused.}
\label{fig:framework}
\end{figure}

\subsection{Edge scorer and fusion}
A 2-layer GCN encoder (hidden dimension $64$, dropout $0.4$) produces node embeddings.
Node-centric and edge-centric edge representations are combined by concatenation in the core comparison; gated and cross-attention fusions are ablations only.
A heteroscedastic decoder outputs $(\mu_e,\log\sigma_e^2)$.
Training minimizes a weighted sum of classification cross-entropy, Gaussian NLL, and ranking hinge loss (Supplementary Information~S9, S13).

\subsection{Exact-budget pruning}
Reported runs keep the top $(1-r)|E|$ edges by descending $\mu_e$, with
$n_{\mathrm{keep}}=|E|-\mathrm{round}(r|E|)$.
Estimated uncertainty $\sigma_e$ is not used for ranking.
Optional constraint-aware pruning exists in code but is disabled in the multi-seed comparison.
